\bookchapter{Глава 0: Привет, прежде чем начать}

Прежде чем отправиться в путешествие, давайте обсудим основные правила, чтобы всё, что будет дальше, было вам понятно. Не спешите — заточка топора не задерживает рубку дров.

\booksection{I. Что эта книга может (и не может) для вас сделать?}

Она поможет вам системно разобраться в том, как пользоваться смартфоном, понять принципы работы его функций, причём независимо от того, пользуетесь ли вы Android или iOS — большинство подходов универсальны. Книга проведёт вас через практические упражнения, чтобы вы перешли от «умею пользоваться» к «понимаю, как это работает».

Это не инструкция к конкретной модели телефона, не новости о технологических новинках и уж тем более не руководство о том, как погрузиться в телефон с головой. Честно говоря, если вы ожидаете узнать «как круче смотреть видео на телефоне», вас, возможно, ждёт разочарование — нас больше волнует, как сделать так, чтобы телефон действительно служил вам, а не наоборот.

\booksection{II. Наши условные обозначения}

Для ясности в книге используются небольшие пометки — это наш с вами «секретный код»:

\begin{itemize}
  \item \textbf{[Название кнопки]}: обозначает конкретную кнопку или название приложения на экране, например [Настройки], [ОК].
  \item \textbf{>}: обозначает последовательность действий, например «Нажмите [Настройки] > [Звук]» означает: сначала нажмите Настройки, затем Звук.
  \item \textbf{Жирный шрифт}: очень важно! Возможно, ключевое понятие или предупреждение безопасности — обязательно обратите внимание.
  \item \emph{Курсив}: дополнительные пояснения, мысли или метафоры, помогающие лучше понять материал.
\end{itemize}

\booksection{III. Android и iOS — расскажем о обоих}

Знаете что? Некоторые вещи называются по-разному, но делают одно и то же. Например: приложение (Android) и App (iOS), настройки (и там и там), файловый менеджер (Android) и Файлы (iOS). Главное — понимать, что это одно и то же, не запутайтесь в терминах.

Кроме того, в лагере Android теперь есть HarmonyOS от Huawei со своими особенностями — мы расскажем о ней отдельно позже. В общем, какую бы марку телефона вы ни использовали, эта книга вам поможет.

\booksection{IV. Объём памяти телефона: почему купил 128 ГБ, а доступно меньше 110 ГБ?}

Единицы измерения памяти всем известны: 1 ТБ = 1024 ГБ, 1 ГБ = 1024 МБ, 1 МБ = 1024 КБ… Но есть подвох: производители накопителей обычно считают по основанию 1000, а система — по основанию 1024, поэтому купленный вами накопитель на «128 ГБ» в реальности «усыхает».

Но у уменьшения доступного объёма памяти телефона есть три основные причины:

\begin{enumerate}
  \item \textbf{Система занимает своё место}: операционная система — как несущие стены дома, она должна существовать. Современные системы легко занимают 10-20 ГБ, и это нормально.
  \item \textbf{Предустановленные производителем приложения}: все те приложения, которые идут в новом телефоне, тоже занимают место — некоторые можно удалить, некоторые нет.
  \item \textbf{Потери на разметку и форматирование}: технические причины — как подкладка рюкзака тоже занимает объём.
\end{enumerate}

\begin{tcbinfo}{Совет при покупке телефона}
Сейчас 128 ГБ уже практически недостаточно, надёжнее начинать с 256 ГБ. Если вы любите фотографировать и играть в игры, советуем сразу брать 512 ГБ или даже больше. Запомните: памяти лучше пусть будет с запасом, чем в обрез.

\end{tcbinfo}

\begin{tcbpractice}{Практика: посмотрите «паспорт» вашего телефона}
Откройте [Настройки], найдите «О телефоне» (или аналогичный пункт) и запишите: модель телефона, версию системы, общий и свободный объём памяти. Это первый шаг к знакомству с ним. Заодно проверьте состояние аккумулятора, чтобы быть в курсе.

\end{tcbpractice}

\booksection{V. Ваш телефон — это «хранилище цифровых активов»}

Прежде чем перейти к конкретным знаниям, давайте усвоим ключевую мысль:\textbf{телефон давно перестал быть просто средством связи — это наше самое важное хранилище цифровых активов}. Контакты, история переписки, семейные фотографии, рабочие документы, банковские приложения, платёжные коды, данные о здоровье, электронные удостоверения… Всё это вместе часто стоит дороже самого телефона.

\begin{tcbinfo}{Почему это важно?}
Если телефон потеряется или сломается, вы потеряете не только «железо» за несколько тысяч рублей,\textbf{но, что гораздо важнее, — накопленные за годы невосполнимые данные}: фотографии роста детей, голосовые сообщения ушедших близких, важные договоры без резервной копии. Грамотная организация и привычка к защите защитят вас в трёх измерениях:

\begin{itemize}
  \item \textbf{Повышение эффективности работы}: упорядоченная классификация файлов ускоряет поиск, телефон долго не тормозит.
  \item \textbf{Продление срока службы}: разумное управление памятью снижает нагрузку на чип, аккумулятор и память служат дольше.
  \item \textbf{Обеспечение безопасности данных}: многоуровневое управление + регулярное резервное копирование + ограничение разрешений — и «катастрофа данных» обходит вас стороной.
\end{itemize}

\end{tcbinfo}

\booksubsection{Наш ключевой принцип: «Профилактика прежде всего, лечение — в дополнение»}

Как нельзя ждать болезни, чтобы пойти к врачу, так и\textbf{управление телефоном нужно начинать с «повседневных привычек»}. В этой книге мы раз за разом подчёркиваем не какие-то продвинутые приёмы, а эти простые принципы:

\begin{itemize}
  \item \textbf{Профилактика прежде всего}: устанавливайте приложения только из официальных магазинов, перед отправкой фото очищайте метаданные, не повторяйте пароли — эти «мелочи» почти ничего не стоят, но блокируют подавляющее большинство рисков.
  \item \textbf{Регулярная проверка}: раз в месяц — уборка папки загрузок, раз в квартал — аудит разрешений, раз в полгода — полное резервное копирование. Фиксированный ритм в сто раз эффективнее, чем авральное тушение пожаров.
  \item \textbf{Всегда имейте запасной выход}: включите облачную синхронизацию, храните учётные данные, составьте «список цифрового наследства» — чтобы в случае неприятностей и вы, и ваши близкие могли спокойно справиться с ситуацией.
\end{itemize}
